\subsection{Problématique}

\begin{itemize}
	\item Sur la compréhension du cinéma par l’IA :  
	
	Ce qui est encore en cours de développement pour les outils actuellement = compréhension du langage cinématographique (montage, narration), analyse des ambiances et émotions, détection d’œuvres musicales. 
	
	Incapacité des machines à comprendre le langage ciné pour l’instant et à l’analyser correctement : pourquoi ?  
	
	Il n’existe pas de solution tout en un  
	
	L'image animée regroupe bcp de choses.
	
	  
	
	\item Interface et l’expérience utilisateur :  
	
	Comment combiner performance technique et expérience utilisateur au sein de l’application avec l’implémentation de l’IA ? 
	
	Implémentation de deux outils de recherche assistée dans l’application : comment guider l’utilisateur vers l’outil qu’il cherche ?  
	
	Faire une interface utilisateur qui augmente la productivité et l’interactivité mais en même temps qui reste adaptée aux institutions culturelles qui n’ont pas forcément d’objectifs de rendement et de productivité comme des entreprises privées. Trouver la limite. Ne pas en faire une appli de marketing.  
	
	Tensions entre besoins utilisateurs musée / et une logique de productivité d’entreprise privée vers laquelle se tourne de plus en plus l’application : tableau de bord, intelligence artificielle.  
	
	Améliorer l’interface et l’expérience utilisateur est ce qui fait prospérer skinsoft et ce qui la démarque de ses concurrents.  
	
	
	\item Migration de données :  
	
	Décalage entre ce que le client veut pour ses données, et ce que skin soft est capable de faire. Un changement de logiciel induit forcément une modification de la saisie des données, même si l’application est entièrement paramétrable. Décalage entre ce que l’institution attend et ce que skinsoft peut faire conformément aux standards documentaires, au temps qu’ils peuvent y consacrer et aux limites de l’application et de l’interface.  
	
	Difficulté de paramétrer l’application en fonction de la manière de saisir les données propres à chaque institution, des thésaurus utilisés... 
	
	L'application et les paramétrages évoluent presque à chaque nouveau client et devient plus malléable.  
	
	
	
	\item Sur l’indexation :  
	
	La génération de résumé sur une archive audiovisuelle prévue par skin soft se fait de quelle façon, à partir de quelles données ? 
	
	Résumé à partir de notices déjà constituées sur l’archives en question 
	
	Extraction des éléments dans l’image : autorités 
	
	Extraction d’éléments dans le son 
	
	Comment le NLP transforme les pratiques d’indexation ? Lien entre les deux ? 
	
	Pour l'instant, il existe des modèles qui permettent de détecter la séparation etre les plans grâce au changement de position des personnages, changement de lumières...
	Mais la détection du changement de séquences se situe à un niveau supplémentaire de compréhension du langage cinématographique. Cela suppose de comprendre ce qu'est une séquence (combinaison d'un changement temporel et spatial). Comment améliorer les modèles qui détecte les changements de plans pour leur faire comprendre les chngement de séquences ? Surtout sur des archives audiovisuelles où il se passe peu de choses dans l'image. La plupart des cas d'usages vus et lus sont entraînés sur des séquences de télévision par exemple, un langage "classique et normé qui ets toujours le même. 
	Pour les archives très anciennes comme celles d'albert kahn, il n'y a pas de coupures ou très peu, les plans sont fixes, les sujets sont généralement loins de la caméra. est ce que c'est plus facile à entraîner ? comment entraîner un modèle qui puisse analyser des vidéos à la fois très similaires et très différentes en termes de langage ciématographique ?
	
	les archives audiovisuelles nécessitent une indexation de contenu : extraction d'information à partir de ce qui setrouve dans l'image
	
	pour les archives papier, l'indexation automatique des archives nées numériques et les archives numérisées est différentes. Pour les arches nées numérique,s il existe déjà un produit numérique : pdf par exemple ou document texte. Pour une archve numérisée, il faut produire ce dérivé pour y appliquer de la reconnaissance d'entités, de caractères... 
	Est-ce la même chose pour les archives audiovisuelles ? Ou ça ne change rien ? 
	
\end{itemize}